\section*{Chapitre \ref{ch:fm:the-multi-manager-platform} --- La plateforme multigérants}

\begin{solution}{exo:fm:the-multi-manager-platform:1}
$0.20\times0.12\times500=12$ millions de dollars.
\end{solution}

\begin{solution}{exo:fm:the-multi-manager-platform:2}
La perte de 25~millions de dollars est d'abord absorbée~: $0.20\times(60-25)=7$ millions de dollars.
\end{solution}

\begin{solution}{exo:fm:the-multi-manager-platform:3}
$\E[\max(X,0)]=\sigma\,(S\Phi(S)+\varphi(S))=\sigma(0.8413+0.2420)=1.0833\,\sigma$.
\end{solution}

\begin{solution}{exo:fm:the-multi-manager-platform:4}
La part de l'équipe gagnante, 2~millions de dollars, plus les 6~millions de dollars de coûts des deux pods~: 8~millions de dollars, pour un profit nul du fonds.
\end{solution}

\begin{solution}{exo:fm:the-multi-manager-platform:5}
$p(\varphi(S)-S\Phi(-S))/S=0.2\times(0.3521-0.5\times0.3085)/0.5=7.9~\%$ du profit brut.
\end{solution}

\begin{solution}{exo:fm:the-multi-manager-platform:6}
Un drawdown de 10~\% représente une demi-année de volatilité pour un pod à 20~\%, atteint dans l'année par la plupart des pods, quel que soit leur edge. Des seuils d'environ une
volatilité, une coupe à 20~\% et un arrêt à 30~\% coupent environ 26~\% des années-équipes dans le modèle.
\end{solution}

\begin{solution}{exo:fm:the-multi-manager-platform:7}
Les rémunérations des équipes passent de 4.27~\$ à 4.04~\$ pour 100~\$, et la part de l'investisseur passe de 9.44~\$ à 9.63~\$~; le coût de compensation tombe à 0.1~\%
du profit brut. Les pertes reportées viennent en déduction des profits ultérieurs de la même équipe, de sorte que rémunérer les gagnants coûte peu lorsque les perdants restent en place.
\end{solution}

\begin{solution}{exo:fm:the-multi-manager-platform:8}
Les historiques de performance sont mesurés sur les survivants et dans l'échantillon~; les équipes qui vont perdre sont celles qui n'ont pas encore été remplacées, et les ratios de Sharpe des nouvelles équipes sont
inconnus. Le risque de compensation dépend du ratio de Sharpe des pods tels qu'ils seront, et de la rotation~: remplacer les perdants le rétablit.
\end{solution}

\begin{solution}{pb:fm:the-multi-manager-platform:1}
\begin{enumerate}
\item Une refacturation des coûts réels du fonds, rémunérations des équipes comprises, aux investisseurs~; la part du profit propre d'une équipe, nette des coûts et des pertes reportées,
versée à l'équipe.
\item Un pourcentage du profit propre de l'équipe l'année précédente, valorisé en mark-to-market, indépendamment des autres équipes et du fonds~; les pertes sont reportées,
en général sans récupération (clawback).
\item Des investissements directs~: couvertures ou transactions de contrepartie contre des expositions partagées, et transactions qui ajoutent aux positions compensées détenues par plusieurs équipes.
\item Chaque équipe est rémunérée sur son propre profit~; les gagnants sont payés pendant que les pertes des perdants retombent sur le fonds.
\item 20.09~\$ de brut~; 4.27~\$ de rémunérations d'équipes~; 4.00~\$ de coûts~; 2.38~\$ pour le gérant~; 9.44~\$ pour l'investisseur.
\item 9 sur 400.
\item 14.47~\$, pour 5.62~\$ de frais~; le gérant paierait 8.27~\$ de rémunérations et de coûts avec 5.62~\$ et perdrait 2.65~\$ pour 100~\$.
\item 47.0~\%.
\item Le coût de la rémunération des équipes sur leur propre profit plutôt que sur le total~;
$p\sigma(\varphi(S)-S\Phi(-S))$ par pod lorsque le fonds est profitable.
\item 1.7~\%, 7.9~\% et 22.9~\% du profit brut.
\item 1.3~\% et 6.2~\%.
\item 0.1~\%~: une équipe qui reste doit regagner ses pertes avant d'être à nouveau rémunérée.
\item La qualité de ses équipes (leurs ratios de Sharpe) et la fréquence à laquelle il remplace les équipes perdantes.
\item Volatilité de 6.9~\% et ratio de Sharpe de 2.9 sans~; 2.7~\% et 7.5 avec~; tout le risque partagé est un seul facteur négociable dans le modèle.
\item 88.4~\% coupés et 26.6~\% arrêtés~; rendement de 20.1~\% à 13.4~\%, volatilité de 6.9~\% à 4.8~\%, ratio de Sharpe de 2.92 à 2.77.
\item Chaque pod a un edge vivant, donc arrêter un pod ne fait que retirer du profit.
\item Le trading propre du centre avec le capital du fonds~; le fonds en est propriétaire, aucune équipe n'est rémunérée dessus.
\item Dans les mêmes limites que les pods, car il peut ajouter à leur concentration.
\item $p\sigma(\varphi(S)-S\Phi(-S))$ par pod, 1.7~\% du profit brut pour $S=1$ et 7.9~\% pour $S=0.5$~; les rémunérations plus les coûts absorbent 41.2~\% du profit brut.
\item La fréquence de remplacement des équipes et leur qualité, qui fixent le coût de compensation~; et ce que fait le \omterm{def:fm:the-multi-manager-platform:centre}{book central} et ce qui le limite.
\end{enumerate}
\end{solution}

\begin{solution}{iq:fm:the-multi-manager-platform:1}
Les coûts du fonds, rémunérations des équipes comprises, refacturés à leur coût aux investisseurs~; le talent et l'infrastructure coûtent plus cher qu'une commission fixe ne peut couvrir.
\iqlookfor{que les frais suivent la base de coûts, et qui la supporte.}
\end{solution}

\begin{solution}{iq:fm:the-multi-manager-platform:2}
Rémunération sur 12~\% moins les 5~\% reportés~: 7~\% du capital multiplié par le taux. Le successeur d'une équipe remplacée démarre sans le carry.
\iqlookfor{le report des pertes et sa remise à zéro.}
\end{solution}

\begin{solution}{iq:fm:the-multi-manager-platform:3}
$\sum_i\max(X_i,0)\ge\max(\sum_i X_i,0)$~; avec un total positif, le coût vaut $p\sum_i(\E\max(X_i,0)-\E X_i)=p\sigma(\varphi(S)-S\Phi(-S))$ par pod.
\iqlookfor{la convexité de la rémunération et la forme fermée.}
\end{solution}

\begin{solution}{iq:fm:the-multi-manager-platform:4}
$0.2\sqrt{(1+49\times0.1)/50}=6.9~\%$~; couvrir la part commune laisse $0.2\sqrt{0.9/50}=2.7~\%$.
\iqlookfor{la formule de variance à corrélation égale.}
\end{solution}

\begin{solution}{iq:fm:the-multi-manager-platform:5}
C'est une demi-année de volatilité, atteinte par la plupart des pods par hasard. On fixe des seuils en unités de volatilité, environ une volatilité, et on agit progressivement.
\iqlookfor{des seuils exprimés en unités de la volatilité propre du pod.}
\end{solution}

\begin{solution}{iq:fm:the-multi-manager-platform:6}
Son P\&L face aux expositions cumulées des pods~: une couverture pure a une corrélation négative avec le P\&L du facteur commun des pods
et une faible dérive propre~; les positions qui grossissent avec le consensus des pods s'y ajoutent.
\iqlookfor{la corrélation avec le book compensé et la taille des positions.}
\end{solution}
